\documentclass[16pt]{article}
\usepackage{amsmath}
\usepackage[utf8]{inputenc}
\usepackage{array}


\begin{document}

\section*{Exercice 01:}

\renewcommand{\arraystretch}{2}

\begin{tabular}{lp{15cm}r}
\hline
\textbf{Formule 1} & On a $12 \div 3 = 4$ signifie aussi $3 \times 4 = 12$.\\ \hline
\textbf{Formule 2} & L'égalité $x = 2y$ est équivalente à $y = x/2$.\\\hline
\textbf{Formule 3} & L'égalité $x = 2y$ est équivalente à $y = \frac{1}{2}x$.\\ \hline
\textbf{Formule 4} & La fonction qui à $x$ associe $x^2$.\\ \hline
\textbf{Formule 5} & On considère la fonction $f$ définie sur \textbf{R}. $$f : x\mapsto e^{2x+1}$$\\ \hline
\textbf{Formule 6} & $$\frac{\frac{x^2}{y^2}}{\frac{y^2}{x^2}} = \frac{x^4}{y^4}$$\\ \hline
\textbf{Formule 7} & $$\dfrac{\dfrac{x^2}{y^2}}{\dfrac{y^2}{x^2}} = \frac{x^4}{y^4}$$\\ \hline
\textbf{Formule 8} &
Dans le repère $(O;\vec{i};\vec{j})$,les sous-espaces $\mathcal{F}$ et $\mathcal{G}$
vérifient $\forall (\vec{u},\vec{v}) \in \mathcal{F} \times \mathcal{G}, \ \vec{u}\cdot\vec{v} = 0.$\\ \hline
\textbf{Formule 9} & $$\frac{\sqrt{x+1}}{y+1} \neq \sqrt{\frac{x+1}{y+1}}$$\\ \hline
\textbf{Formule 10}& $$1 + \frac{1}{1 + \dfrac{1}{1 + \dfrac{1}{1 + \dfrac{1}{1}}}} = 1,6$$\\ \hline
\end{tabular}

\section*{Exercice 02:}
\begin{tabular}{lp{15cm}}
\textbf{Formule 1} & On a $\sin \frac{\pi}{2} = 1$ et $\sin \frac{\pi}{4} = \frac{\sqrt{2}}{2}$ \\ \hline
\textbf{Formule 2} &
\[ \sum_{i,j\in I\times J} i+j = \sum_{i\in I}\left(\sum_{j\in J} i+j\right) \]\\ \hline
\textbf{Formule 3} &
\[ \sum_{p\in\mathcal{P}}\frac{1}{p} \qquad\text{ou}\qquad
   \sum_{p\text{ premier}}\frac{1}{p} \]\\ \hline
\textbf{Formule 4} \quad &
Par définition, \(nx=\overbrace{x+\cdots+x}^{n\text{ fois}}\).\\ \hline
\textbf{Formule 5} \quad &
On a : \(nx \overset{\text{déf}}{=} \overbrace{x+\cdots+x}^{n\text{ fois}}\). \\ \hline
\medskip
\textbf{Formule 6} & \[ f : \begin{array}{c} \mathbf{R}\to\mathbf{R} \\ x\mapsto e^x \end{array} \]
On verra lors de la correction comment obtenir le résultat(plus classique) :\[ f : \begin{array}{rcl} \mathbf{R} & \to & \mathbf{R} \\
                          x & \mapsto & e^x \end{array} \]\\ \hline
\textbf{Formule 7} &
\[ \iint_{\mathcal{D}} z^2\,\mathrm{d}z \]
La difficulté, ici, est de gèrer correctement les espacements.\\ \hline
\textbf{Formule 8} &
\[ \underbrace{\overbrace{a,\ldots,a}^{k},\overbrace{c,\ldots,c}^{l}}_{k+l\text{ éléments}} \]\\ \hline
\textbf{Formule 9} &
\[
\begin{pmatrix}
  \begin{pmatrix} 1 & 0 \\ 0 & 1 \end{pmatrix} &
  \begin{pmatrix} a & b \\ c & d \end{pmatrix} \\
  \begin{pmatrix} d & -b \\ -c & a \end{pmatrix} &
  \begin{pmatrix} 0 & -1 \\ 1 & 0 \end{pmatrix}
\end{pmatrix}
\]\\ \hline
\end{tabular}


\section*{Exercice 03:}
\begin{center}
\begin{tabular}{lp{4cm}rr}
\textbf{Corps} & \textbf{Commentaire} & \textbf{Ø (km)} & \textbf{densité} \\
Soleil & Il s'agit d'une étoile tout à fait quelconque. & 1392000 & 1,409 \\
Mercure & Il fait très chaud. & 4840 & 5,50 \\
Vénus & Souvent appelée " jumelle " de la terre. C'est l'étoile du berger. & 12390 & 5,25 \\
Terre & Toutes les personnes de ce stage l'habite. & 12760 & 5,517 \\
Mars & Peut-être la prochaine planète colonisée par l'homme ? & 6800 & 3,94 \\
\end{tabular}
\end{center}



\begin{center}
\begin{tabular}{|l|>{\raggedright\arraybackslash}p{4cm}|r|r|}
\hline
\textbf{Corps} & \textbf{Commentaire} & \textbf{Ø (km)} & \textbf{densité} \\
\hline \hline
Soleil & Il s'agit d'une étoile tout à fait quelconque. & 1392000 & 1,409 \\
\hline
Mercure & Il fait très chaud. & 4840 & 5,50 \\
Vénus & Souvent appelée " jumelle " de la terre. C'est l'étoile du berger. & 12390 & 5,25 \\
Terre & Toutes les personnes de ce stage l'habite. & 12760 & 5,517 \\
Mars & Peut-être la prochaine planète colonisée par l'homme ? & 6800 & 3,94 \\
\hline
\end{tabular}
\end{center}


\end{document}


